%------------------------------------------------------------------------------%
\documentclass[a4paper,12pt,fleqn]{article}

\usepackage{calculo_varias_variaveis-1}

\ShowAnswers

%------------------------------------------------------------------------------%
\begin{document}

\makeHeader{CFVVI}{Exame Especial -- Turma 2}{28/07/2025}

% 01
%------------------------------------------------------------------------------%
\exercise{25}
Considerando que a equação
\(
  xe^y + ye^z + 2\ln(x) -2 -3\ln(2) = 0
\)
define $z$ como função de $x$ e $y$.
\begin{enumerate}[label=\alph*)]
  \item Calcule as derivadas \(\D{z}{x}\) e \(\D{z}{y}\)
  \item Encontre os valores de \(\D{z}{x}\) e \(\D{z}{y}\) em \((1, \ln(2), \ln(3))\)
  \item Construa a aproximação linear da função $z(x, y)$ no ponto \((1, \ln(2))\)
\end{enumerate}
\clearpagequestiononly

\begin{answer}
  \vspace{\baselineskip}
  \textbf{a)\;}
  Considerando $z=z(x, y)$ e derivando por $x$ os dois lados da equação temos
  \begin{align*}
    xe^y + ye^{z(x, y)} + 2\ln(x) & = 2 +3\ln(2)  \\
    \D{}{x}\left[xe^y + ye^{z(x, y)} + 2\ln(x)\right] & = \D{}{x}\left[2 +3\ln(2)\right] \\
    \D{}{x}\left(xe^y\right) +
    \D{}{x}\left(ye^{z(x, y)}\right) +
    \D{}{x}\left(2\ln(x)\right)      & = 0 \\
    e^y + ye^{z(x, y)}\D{z}{x} + \frac{2}{x} & = 0 \\
    ye^{z(x, y)}\D{z}{x} & = -e^y - \frac{2}{x}  \\
    \D{z}{x}     & = \frac{-1}{ye^{z(x, y)}}\left(e^y + \frac{2}{x}\right)
  \end{align*}
  Derivando por $y$ os dois lados da equação temos
  \begin{align*}
    xe^y + ye^{z(x, y)} + 2\ln(x) & = 2 +3\ln(2) \\
    \D{}{y}\left[xe^y + ye^{z(x, y)} + 2\ln(x)\right] & = \D{}{y}\left[2 +3\ln(2)\right] \\
    \D{}{y}\left(xe^y\right) +
    \D{}{y}\left(ye^{z(x, y)}\right) +
    \D{}{y}\left(2\ln(x)\right) & = 0 \\
    x\D{e^y}{y} + \D{y}{y}e^{z(x, y)} + y\D{}{y}e^{z(x, y)} + 0 & = 0 \\
    xe^y + e^{z(x, y)} + ye^{z(x, y)}\D{z}{y} & = 0 \\
    ye^{z(x, y)}\D{z}{y} & = -xe^y - e^{z(x, y)} \\
    \D{z}{y} & = - \frac{xe^y + e^{z(x, y)}}{ye^{z(x, y)}}
  \end{align*}

  \vspace{\baselineskip}
  \textbf{b)\;}
  Avaliando as derivadas no ponto \((1, \ln(2), \ln(3))\)
  \begin{align*}
    \D{z}{x}(1, \ln(2))
    & = \left(\frac{-1}{ye^z}\left(e^y + \frac{2}{x}\right)\right)\evalat{(1, \ln(2), \ln(3))}{} \\
    & = \frac{-1}{\ln(2)e^{\ln(3)}}\left(e^{\ln(2)} + \frac{2}{1}\right) \\
    & = \frac{-1}{3\ln(2)}\left(2 + 2\right) \\
    & = \frac{-4}{3\ln(2)} \\
    \D{z}{y}(1, \ln(2))
    & = \left(- \frac{xe^y + e^{z(x, y)}}{ye^{z(x, y)}}\right)\evalat{(1, \ln(2), \ln(3))}{} \\
    & = - \frac{1e^{\ln(2)} + e^{\ln(3)}}{{\ln(2)}e^{\ln(3)}} \\
    & = - \frac{2 + 3}{{\ln(2)}3} \\
    & = \frac{-5}{3\ln(2)}
  \end{align*}

  \vspace{\baselineskip}
  \textbf{c)\;}
  A aproximação linear da função $z(x, y)$ no ponto \((x_0, y_0) = (1, \ln(2))\) é
  \begin{align*}
    L(x, y)
    & = z(x_0, y_0) + (x-x_0)\D{z}{x}(x_0, y_0) + (y-y_0)\D{z}{y}(x_0, y_0) \\
    & = z(1, \ln(2)) + (x-1)\D{z}{x}(1, \ln(2)) + (y-\ln(2))\D{z}{y}(1, \ln(2)) \\
    & = \ln(3) + (x-1)\frac{-4}{3\ln(2)} + (y-\ln(2))\frac{-5}{3\ln(2)} \\
    & = \ln(3) - \frac{4}{3\ln(2)}(x-1) - \frac{5}{3\ln(2)}(y-\ln(2))
  \end{align*}

  Não faz parte da resolução solicitada na prova, mas com o uso de uma calculadora
  podemos avaliar essa expressão, obtendo uma aproximação para $z(x, y)$
  quando $(x, y)\approx (1, 0.69)$
  \begin{align*}
    z(x, y) & \approx L(x, y) \\
    & = \ln(3)
    + \frac{4}{3\ln(2)}
    + \frac{5}{3}
    - \frac{4}{3\ln(2)} x
    - \frac{5}{3\ln(2)} y \\
    & \approx 4,689 -1,924 x - 2,404 y
  \end{align*}
\end{answer}

% 02
%------------------------------------------------------------------------------%
\exercise{25}
Encontre os valores máximo e mínimo da função
\(
  f(x, y) = 2x^2 + y^2
\)
restrita a região fechada limitada
\(
  x^2 + 2y^2 \leq 16
\)
\clearpagequestiononly

\begin{answer}
  \setlength\columnsep{3em}
  \setlength{\jot}{5pt}
  \begin{multicols}{2}
    \raggedcolumns
    Como a função é contínua e a região é fechada e limitada sabemos que
    $f$ assume um valor máximo e um valor mínimo na região.

    \vspace{\baselineskip}
    Temos que buscar os pontos críticos no interior e aplicar multiplicadores
    de Lagrange na fronteira.

    \vspace{\baselineskip}
    \textbf{Gradiente} de $f$
    \[
      \D{f}{x}
      = \D{}{x}\left(2x^2 + y^2\right)
      = 4x
    \]
    \[
      \D{f}{y}
      = \D{}{y}\left(2x^2 + y^2\right)
      = 2y
    \]

    \vspace{\baselineskip}
    \textbf{Pontos críticos}:
    Impondo \(\nabla f = 0\) temos o ponto interior
    \[
      P_1 = (0, 0)
    \]

    \vspace{\baselineskip}
    \textbf{Multiplicadores de Lagrange}:
    Temos a restrição \(g(x, y) = x^2 + 2y^2 \leq 16\) e seu gradiente é
    \[
      \D{g}{x}
      = \D{}{x}\left(x^2 + 2y^2\right)
      = 2x
    \]
    \[
      \D{g}{y}
      = \D{}{y}\left(x^2 + 2y^2\right)
      = 4y
    \]

    Precisamos resolver o sistema
    \begin{gather*}
      4x = 2\lambda x \\
      2y = 4\lambda y \\
      x^2 + 2y^2 = 16
    \end{gather*}
    Que pode ser simplificado
    \begin{gather*}
      2x =  \lambda x \\
      y = 2\lambda y \\
      x^2 + 2y^2 = 16
    \end{gather*}
    Da primeira equação temos que $x=0$ ou $\lambda=2$.

    Se $x=0$ a terceira equação se reduz a
    \begin{align*}
      x^2 + 2y^2 & = 16           &  \\
      2y^2       & = 16           &  \\
      y^2        & = 8            &  \\
      y          & = \pm 2\sqrt{2} &
    \end{align*}
    Substituindo qualquer uma delas na segunda equação
    temos $\lambda = 0$. Encontramos o segundo e terceiro pontos
    \[
      P_2 = \left(0, -2\sqrt{2}\right)
      \qquad
      P_3 = \left(0, 2\sqrt{2}\right)
    \]

    Se $\lambda=2$ a segunda equação se torna $y=8y$ e, portanto, $y=0$.
    Substituindo na terceira equação temos $x^2=16$, cujas soluções são
    $x=-4$ ou $x=4$.  Encontramos o quarto e quinto pontos
    \[
      P_4 = (-4, 0)
      \qquad
      P_5 = (4, 0)
    \]

    \textbf{Avaliando} a função nos pontos encontrados
    \begin{align*}
      f( 0,  0)                    & = 2\times    0^2 +    0^2 =  0                 \\
      f\left( 0, -2\sqrt{2}\right) & = 2\times 0^2 + \left(-2\sqrt{2}\right)^2 = 8 \\
      f\left( 0,  2\sqrt{2}\right) & = 2\times 0^2 + \left( 2\sqrt{2}\right)^2 = 8 \\
      f(-4,  0)                    & = 2\times (-4)^2 +    0^2 = 32                 \\
      f( 4,  0)                    & = 2\times    4^2 +    0^2 = 32
    \end{align*}

    O valor mínimo de $f$ é 0 e o valor máximo é 32
  \end{multicols}
\end{answer}

% 3
%------------------------------------------------------------------------------%
\exercise{25}
Considerando a função
\(
  f(x, y) = \sqrt{9-x^2-y^2}-3
\)
\begin{enumerate}[label=\alph*)]
  \item Determine e esboce o domínio de $f$
  \item Determine as curvas de nível de $f$ e esboce três delas
  \item Calcule o limite
  \(
    \lim_{(x, y)\to(0, 0)} \frac{f(x, y)}{x^2+y^2}
  \)
  ou mostre que ele não existe
\end{enumerate}
\begin{questiononly}
  \begin{multicols}{2}
    Domínio
    \vspace{-0.8\baselineskip}
    \begin{center}
      \includegraphics{axis_xy}
    \end{center}
    Curvas de nível
    \vspace{-0.8\baselineskip}
    \begin{center}
      \includegraphics{axis_xy}
    \end{center}
  \end{multicols}
\end{questiononly}
\clearpagequestiononly

\begin{answer}
  \begin{multicols}{2}
    Domínio
    \vspace{-0.8\baselineskip}
    \begin{center}
      \includegraphics{T4_dominio}
    \end{center}
    Curvas de nível
    \vspace{-0.8\baselineskip}
    \begin{center}
      \includegraphics{T4_curvas_nivel}
    \end{center}
  \end{multicols}

  \setlength\columnsep{3em}
  \begin{multicols}{2}
    \raggedcolumns
    \setlength{\jot}{5pt}
    \vspace{\baselineskip}
    \textbf{a)\;}
    O domínio de $f$ consiste dos pontos onde é possível avaliar
    a raiz quadrada, isto é, os pontos que satisfazem $9-x^2-y^2\geq 0$,
    ou seja,
    \begin{align*}
      9 -x^2 -y^2 & \geq 0   \\
      -x^2 -y^2   & \geq -9 \\
      x^2 +y^2    & \leq 9 = 3^2
    \end{align*}
    Portanto,
    \[
      D_f = \left\{(x, y)\in\R^2 \st x^2 + y^2 \leq 9 \right\},
    \]
    que corresponde a um disco de raio 3, centrado na origem.

    \vspace{\baselineskip}
    \textbf{b)\;}
    As curvas de nível de $f$ são compostas pelos pontos onde $f(x, y) = c$
    para alguma constante $c$ na imagem de $f$.

    Para determinar a imagem de $f$ observamos que $x^2+y^2$ só pode assumir
    valores no intervalo $[0, 9]$. Portanto, $9 - x^2 - y^2$ está limitado
    ao mesmo intervalo. Consequentemente $\sqrt{9 - x^2 - y^2}$ está em $[0, 3]$.
    Concluímos que os valores de $f$ estão em $[-3, 0]$, isto é,
    \(\Im(f) = [-3, 0]\)

    Para qualquer $c\in\Im(f)=[-3, 0]$, temos a curva de nível
    \begin{align*}
      f(x, y)                & = c             \\
      \sqrt{9 -x^2 -y^2} - 3 & = c             \\
      \sqrt{9 -x^2 -y^2}     & = c + 3         \\
      9 -x^2 -y^2            & = (c + 3)^2     \\
      -x^2 -y^2              & = (c + 3)^2 -9  \\
      x^2 +y^2               & = 9 - (c + 3)^2
    \end{align*}
    Que corresponde a uma circunferência centrada na origem de raio
    \[
      r = \sqrt{9 - (c + 3)^2}
    \]

    Escolhendo os valores $c=-3$, $c=-2$ e $c=-1$, temos as curvas de nível
    \begin{align*}
      \gamma_1\colon & x^2 + y^2 = 9 - (-3 + 3)^2 = 9 - 0 = 9 \\
      \gamma_2\colon & x^2 + y^2 = 9 - (-2 + 3)^2 = 9 - 1 = 8 \\
      \gamma_3\colon & x^2 + y^2 = 9 - (-1 + 3)^2 = 9 - 4 = 5
    \end{align*}

    \vspace{\baselineskip}
    \textbf{c)\;}
    Calculando o limite
    \setlength{\jot}{12pt}
    \setlength{\mathindent}{0cm}
    \begin{align*}
      L
      & = \lim_{(x, y)\to(0, 0)} \frac{f(x, y)}{x^2+y^2} \\
      & = \lim_{(x, y)\to(0, 0)} \frac{\sqrt{9-x^2-y^2}-3}{x^2+y^2} \\
      & = \lim_{(x, y)\to(0, 0)} \frac{\sqrt{9-x^2-y^2}-3}{x^2+y^2} \times \\
      & \hspace{32mm}            \frac{\sqrt{9-x^2-y^2}+3}{\sqrt{9-x^2-y^2}+3} \\
      & = \lim_{(x, y)\to(0, 0)} \frac{\left(\sqrt{9-x^2-y^2}\right)^2-3^2}
      {\left(x^2+y^2\right)\left(\sqrt{9-x^2-y^2}+3\right)} \\
      & = \lim_{(x, y)\to(0, 0)} \frac{9-x^2-y^2-9}{\left(x^2+y^2\right)\left(\sqrt{9-x^2-y^2}+3\right)} \\
      & = \lim_{(x, y)\to(0, 0)} \frac{-\left(x^2+y^2\right)}{\left(x^2+y^2\right)\left(\sqrt{9-x^2-y^2}+3\right)}\\
      & = \lim_{(x, y)\to(0, 0)} \frac{-1}{\sqrt{9-x^2-y^2}+3} \\
      & = \frac{-1}{\sqrt{9-0^2-0^2}+3} \\
      & = \frac{-1}{3 + 3} = \frac{-1}{6}
    \end{align*}
  \end{multicols}
\end{answer}

% 04
%------------------------------------------------------------------------------%
\exercise{25}
Dado \(z=\left(\sqrt{3}+i\right)^4 \), calcule
\begin{enumerate}[label=\alph*),itemsep=0pt]
  \item a parte real de $z$,
  \item a parte imaginária de $z$,
  \item o módulo de $z$,
  \item o argumento de $z$
\end{enumerate}
\clearpagequestiononly

\begin{answer}
\setlength\columnsep{3em}
\begin{multicols}{2}
  \raggedcolumns
  \setlength{\jot}{5pt}
  Convertendo $u=\sqrt{3}+i$ para a forma polar
  \[
    \rho
    = \abs{u}
    = \sqrt{\left(\sqrt{3}\right)^2 + 1^2}
    = \sqrt{4}
    = 2
  \]
  \[
    \sin(\varphi) = \frac{\Im(u)}{\abs{u}} = \frac{1}{2}
  \]
  \[
    \cos(\varphi) = \frac{\Re(u)}{\abs{u}} = \frac{\sqrt{3}}{2}
  \]
  \[
    \varphi = \arg(u) = \ang{30} = \frac{\pi}{6}
  \]
  Assim
  \begin{align*}
    u
    & = \rho\left[\cos\left(\varphi\right)+ i\sin\left(\varphi\right)\right] \\
    & = 2\left[\cos\left(\frac{\pi}{6}\right)+ i\sin\left(\frac{\pi}{6}\right)\right]
  \end{align*}
  Portanto
  \begin{align*}
    z
    & = u^3 \\[2mm]
    & = \rho^4\left[\cos\left(4\varphi\right)+ i\sin\left(4\varphi\right)\right] \\[2mm]
    & = 2^4\left[\cos\left(\frac{4\pi}{6}\right)+ i\sin\left(\frac{4\pi}{6}\right)\right] \\[2mm]
    & = 16\left[\cos\left(\frac{2\pi}{3}\right)+ i\sin\left(\frac{2\pi}{3}\right)\right] \\[2mm]
    & = 16\left[-\frac{1}{2}+ i\frac{\sqrt{3}}{2}\right] \\[2mm]
    & = -8 + 8\sqrt{3}i
  \end{align*}
  Parte real de $z$
  \[
    \Re(z) = -8
  \]
  Parte imaginária de $z$
  \[
    \Im(z) = 8\sqrt{3}
  \]
  Módulo de $z$
  \[
    \abs{z} = 16
  \]
  Argumento de $z$
  \[
    \arg(z) = \frac{2\pi}{3}
  \]
\end{multicols}
\end{answer}

%------------------------------------------------------------------------------%
\end{document}
%------------------------------------------------------------------------------%
